\chapter*{Conclusions et perspectives}
\addcontentsline{toc}{chapter}{Conclusion générale} 

et leur application potentielle en électromagnétisme pour l’\efie (Electric Field Integration Equation).

Les propriétés intrinsèques des opérateurs \ddfv, que ce soit en terme d’identités vérifiées au niveau discret ou bien de dimension des espaces présentent un intérêt notamment dans les méthodes de préconditionnement utilisées en \efie.


Lors de la simulation numérique du problème de diffraction électromagnétique par un objet, on est amené à devoir résoudre les équations de Maxwell en régime harmonique à l'extérieur de l'objet considéré. On peut alors utiliser la méthode dite des équations intégrales qui revient, dans sa formulation la plus simple (appelée \efie), à résoudre sur la surface $\surflisse$ de l'objet diffractant
\begin{equation} \label{eq:efie}
\gamma^\mathrm{tan} \mathopen{}\left( \mathrm{i} \, \omega \mu_0 G \vect{J} + \frac{\mathrm{i}}{\omega \varepsilon_0} \grad G \, \mathrm{div} \vect{J} \right) = E^\mathrm{tan}_\mathrm{inc}.
\end{equation}
Dans cette équation, $E^\mathrm{tan}_\mathrm{inc}$ est la composante tangentielle du champ électrique incident, $\gamma^\mathrm{tan}$ est la trace tangentielle, $\omega$ est la pulsation de l'onde incidente, et $\varepsilon_0$ et $\mu_0$ sont les coefficients diélectriques du milieu. Enfin, $\vect{J}$ est le courant électrique (c'est un champ de vecteurs tangents sur la surface $\surflisse$ de l'objet) et $G$ est la convolution sur $\surflisse$ par $\displaystyle\frac{\exp(-\i k \abs{x})}{4 \pi \abs{x}}$ où le nombre d'onde $k$ est donné par $k = \omega / c$ avec $c$ la vitesse de la lumière. 

Toutefois, la résolution numérique de l'équation linéaire \ref{eq:efie} est délicate car l'opérateur sous-jacent
\[
\gamma^\mathrm{tan} T = \gamma^\mathrm{tan} \mathopen{}\left( \i \omega \mu_0 G + \frac{\i}{\omega \varepsilon_0} \grad G \div \right)
\] 
aboutit après discrétisation à des matrices très mal conditionnées. Une méthode originale de préconditionnement À COMPLETER consiste à utiliser $\vect{n} \times T$ où $\vect{n}$ est la normale extérieure à l'objet et qui grâce à l'identité de Calder\'on est son propre inverse à perturbation compacte près. 

Cette propriété provient principalement de la relation continu
\begin{equation} \label{eq:eq2}
\div(\vect{n} \times \grad \varphi) = 0
\end{equation}
pour tout fonction $\varphi$ scalaire, définie sur $\surflisse$, qui est l'analogue surfacique de $\div(\rot \psi) = 0$ pour les champs tridimensionnels. 

La méthode de discrétisation doit donc à tout le moins respecter deux propriétés importantes:
\begin{itemize}
    \item Avoir un analogue de \ref{eq:eq2} au niveau discret.
    \item Faire en sorte que $\vect{u} \mapsto \vect{n} \times \vect{u}$ où $\vect{u}$ est un champ de vecteurs sur $\surflisse$ soit inversible dans l'espace discret (afin que $\vect{n} \times T$ ait une chance d'être inversible aussi). Ce point qui paraît anodin est en fait un des principaux obstacles À COMPLETER.
\end{itemize}

De nombreuses méthodes de discrétisation ont été introduites pour résoudre ce problème parmi lesquelles on peut citer les méthodes de type Nédélec, basées sur des discrétisations de Whitney valables aussi sur des surfaces. Toutefois, ces méthodes ne permettent pas d'utiliser l'identité de Calder\'on, et des adaptations complexes et peu naturelles sont nécessaires pour pouvoir le faire. Un des objectifs de la thèse est ainsi de simplifier ce genre d'approches. 

Les méthodes \ddfv constituent un candidat assez naturel et adapté dans le cadre de l'\efie. En effet, la conservation des propriétés des opérateurs au niveau discret est un élément central des méthodes \ddfv ce qui permet de facilement envisager la conservation de \ref{eq:eq2}. De plus, la prise en compte, directement dans la définition des maillages, de la dualité entre centres et sommets des mailles permet d'obtenir des dimensions d'espaces de discrétisation qui permettent d'envisager l'inversibilité de la version discrète de l'opérateur $\vect{u} \mapsto \vect{n} \times \vect{u}$. 